\chapter{Кредитный риск и роль искусственного интеллекта в его оценке}

%       	
%       Впервые искусственный интеллект возник в 1950 г. в работе  американского ученого А. Тьюринга\footnote{А. Тьюринг (1912-1954 гг.) -- английский математик, логик, криптограф, основоположник ИИ.} «Вычислительные машины и разум», в которой представил тест Тьюринга для оценки способности машины имитировать человеческий разум. Традиционно ИИ как исследовательская область зародился в 1956 г. на летнем семинаре в Дартмутском колледже, который организовали четверо американских ученых Д. Мак-Карти\footnote{Д. Мак-Карти (1927-2011 гг.) -- американский ученый, автор термина «искусственный интеллект», лауреат премии Тьюринга за исследования в области ИИ.}, М. Мински\footnote{М. Мински (1927-2016 гг.) -- американский учёный в области искусственного интеллекта, сооснователь Лаборатории искусственного интеллекта в Массачусетском технологическом институте}, Н. Рочестер\footnote{Н. Рочестер (1919-2001 гг.) -- главный архитектор первого научного компьютера компании IBM, основоположник области ИИ.} и К. Шеннон\footnote{К. Шеннон (1916--2001 гг.) -- американский инженер, криптоаналитик, математик, основатель теории информации.}. В основе лежала идея о том, что все когнитивные способности, как обучение, мышление, восприятие, память и научные открытия могут быть запрограммированы так, что компьютер может самостоятельно воспроизводить эти функции. За более чем 60 лет развития ИИ ученые так и не получили окончательного доказательства, что эту идею можно реализовать, но также и не доказали, что она невозможна.
%       
%       В 1958 г. Г. Саймон\footnote{Г. Саймон (1916-2001 гг.) -- американский учёный в области социальных, политических и экономических наук.} заявил, что с развитием исследований машина сможет одержать победу над человеком в чемпионате мира по шахматам. Однако в 1965 г. десятилетний мальчик одержал победу над компьютером в шахматной партии, но исследования в данной области не приостановились, а начали развиваться в других направлениях. Ученые заинтересовались психологией, механизмами понимания, которые привели к созданию экспертных систем, содержащих знания квалифицированных специалистов. Такие технические усовершенствования привели к разработке алгоритмов машинного обучения (далее МО), с помощью которых компьютеры стали накапливать знания и обучаться на основе предыдущего опыта. Однако такие машины требовали больших вычислительных мощностей и финансирование приостановилось.
%       
%       С конца 1990 г. стали совмещать ИИ и робототехнику с целью создания «умных агентов» с эмоциями и чувствами, что привело к новому исследовательскому направлению эмоциональных вычислений (affective computing), анализирующих эмоциональное состояние собеседника, что позволило улучшить диалоговые системы. Также в 1997 году компания IBM Deep Blue одержала победу над чемпионом мира по шахматам Гарри Каспаровым, что стало важным достижением в продвижении ИИ. 
%       
%       С 2010 года мощность компьютеров существенно возросла и появилось новое направление больших данных (Big Data), приведшее к появлению методов глубокого обучения (далее ГО), подразумевающие использование нейронных сетей.
%       
%       В 2022 году был разработан продвинутый чат-бот ChatGPT\footnote{Generative Pre-trained Transformer -- ИИ, который способен решать различные задачи, предварительно обученный на миллионах данных} от компании OpenAI, использующий ИИ и обработку естественного языка (NLP) для ведения человекоподобных диалогов, выполнения таких задач как перевод текста на другие языки, получение инструкций по выполнению различных задач, сочинение стихов и музки и т.д. Создание такой уникальной технологии подогрело интерес общества к дальнейшему изучению ИИ.

\section{Кредитный скоринг как инструмент управления кредитным риском в банках второго уровня}

Кредитный риск, возникающий вследствие неисполнения или ненадлежащего исполнения заемщиком обязательств перед банком, является одним из ключевых факторов, определяющих устойчивость банковской деятельности. Управление кредитным риском представляет собой комплекс мер по выявлению, оценке, мониторингу клиентов, а также минимизации потенциальных потерь. В данном процессе ключевую роль играет анализ платежеспособности заемщика, насколько своевременно он способен исполнить взятые долговые обязательства. 

Поэтому главным инструментом такой оценки является кредитный скоринг. В научной литературе и современных практиках существует множество определений данного понятия от традиционного как статистической модели до современного подхода как комплексной системы принятия решений на основе ИИ. Систематизация основных определений представлена в таблице \ref{tab:credit_scoring_definitions}.

Анализ данных определений показывает (см. Таблицу \ref{tab:credit_scoring_definitions}), что общим для всех подходов является понимание кредитного скоринга как процесса оценки кредитного риска и инструмента принятия кредитных рещений на основе финансовых данных из прошлого. Новизна современных исследований, ориентированных на применение ИИ, заключается в переходе от статистических методов на методы МО и ГО, а также в расширении источников данных за счет альтернативных признаков таких как коммунальные платежи, оплата за аренду жилья и др.

\newpage
\begin{center}
\begin{longtable}{|p{0.63\textwidth}|p{0.32\textwidth}|}
	\caption{Определение кредитного скоринга в современных практиках}
	\label{tab:credit_scoring_definitions} \\
	\hline
	\textbf{Определение} & \textbf{Исследователи} \\ \hline
	\endfirsthead
	
	\hline
	\textbf{Определение} & \textbf{Исследователи} \\ \hline
	\endhead
	
	Систематический процесс, оценивающий кредитный риск человека на основе различных финансовых факторов и его поведения из прошлого. 
	& Addy W.A., Ajayi-Nifise A.O., Bello B.G., Tula S.T., Odeyemi O., Falaiye T.\footnote{см. Там же. С. 7.} \\ \hline
	
	Фундаментальный инструмент финансового принятия решений, который даёт кредитору ключевое понимание риска потенциального заёмщика. 
	& Bari H.\footnote{см. Там же. С. 8.} \\ \hline
	
	Использование искусственного интеллекта и алгоритмов машинного обучения для оценки кредитоспособности на основе нетрадиционных источников данных. 
	& Omotosho M.A.\footnote{см. Там же. С. 12.} \\ \hline
	
	Процесс количественной оценки кредитоспособности заёмщика, при котором прогноз кредитного риска формируется на основе данных о заёмщике с использованием моделей машинного обучения. 
	& Vakrani D.S., Padhye P.S., Gupta S.K., Roy J.K., Nerlekar V.S., Parashar N.\footnote{см. Там же. С. 13.} \\ \hline
	
	Оценка кредитоспособности на основе моделей AI/ML, которые используют расширенный набор данных, включая альтернативные показатели (например, платежи за аренду и коммунальные услуги), чтобы повысить полноту и справедливость оценки. 
	& Nuka T.F., Ogunola A.A.\footnote{см. Там же. С. 11.} \\ \hline
	
	Это процесс оценки кредитного риска без ручного вмешательства, включающий использование машинного обучения (МО) и аналитики больших данных для анализа огромного массива информации. 
	& Emagia (американская IT-компания)\footnote{\cite{Emagia} (дата обращения: 15.02.2026). Загл. с экрана.} \\ \hline
	
	Статистический анализ, проводимый финансовыми учреждениями для определения кредитоспособности физических и юридических лиц. 
	& Invento Labs (белорусская IT-компания)\footnote{\cite{InventoLabs} (дата обращения: 15.02.2026). Загл. с экрана.} \\ \hline
\end{longtable}
 \begin{minipage}{\textwidth}
	\fontsize{9}{11}\selectfont
	\justifying
	Источник: составлено автором на основе анализа научных исследований и открытых источников.
\end{minipage}
\end{center}

На основе анализа современных подходов к кредитному скорингу, включая внедрение технологий ИИ, в рамках данного исследования сформулировано следующее определение:

Кредитный скоринг -- систематический процесс, позволяющий прогнозировать вероятность дефолта заемщика на основе моделей машинного обучения (далее МО) и глубокого обучения (далее ГО) с использованием традиционных и альтернативных данных, применямый банками второго уровня при принятии кредитных решений.

В настоящее время развитие и внедрение ИИ в Казахстане находятся на этапе становления, что отражается в рейтинге готовности по применению ИИ среди стран. Согласно данным Oxford Insights в 2023 г. страна занимала 72 место\footnote{\cite{Cninform} (дата обращения: 15.02.2026). Загл. с экрана.}, в 2024 г. -- 76 место\footnote{\cite{Oxford} (дата обращения: 15.02.2026). Загл. с экрана.}, в 2025 г. -- 58 место\footnote{\cite{Oxford2} (дата обращения: 15.02.2026). Загл. с экрана.}. В региональном разрезе на 2025 г. республика занимает 3 место\footnote{см. Там же. С. 10.} в Южной и Центральной Азии, уступая место Турции и Индии. Наиболее сильные позиции по внедрению ИИ Казахстан показывает в госсекторе, тогда как практическое внедрение и распространение технологий на основе ИИ в экономике остается низким\footnote{см. Там же. С. 10.}. 

Институциональная база развития ИИ также находится на этапе развития. В 2024 г. была утверждена Концепция развития ИИ на 2024-2029 гг., а в начале 2026 г. вступил в силу закон «Об искусственном интеллекте». Стоит отметить, что в действующем законе\footnote{Об искусственном интеллекте [Электронный ресурс] : [закон Республики Казахстан от 17 ноября 2025 года № 230-VIII : по состоянию на 6.03.2026 г.]. URL: \url{https://adilet.zan.kz/rus/docs/Z2500000230} (дата обращения: 10.03.2026).} отсутствуют нормы, регулирующие применение технологий ИИ в финансовом секторе, в т.ч. в кредитном скоринге. Закон носит общий характер по использованию ИИ независимо от отрасли их применения. Следовательно, на данный момент в Республике Казахстан отсутствует специализированное регулирование ИИ в банковском кредитовании. Однако можно выделить следующие нормы, которые косвенно затрагивают применение ИИ в кредитном скоринге:
\begin{enumerate}[nosep]
	\item Классификация ИИ-систем по степени риска (ст. 17) -- кредитный скоринг, инструмент влияющий на доступ граждан к финансовым ресурсам, относится к системам высокой степени риска, поскольку для таких систем обязателен человеческий контроль и машина не может принимать окончательное решение без подтверждения человека.
	\item Справедливость и равенство (ст. 6) --  системы ИИ должны исключать любую дискриминацию при принятии решений.
	\item Требования к прозрачности и объяснимости (ст. 7) -- заемщик имеет право знать, на основании каких данных и каким образом было принято решение о выдаче кредита.
	\item Ответственность и подконтрольность (ст. 8) -- банк, как владелец ИИ-системы несет ответственность за результаты ее работы и обязан осуществлять регулярный контроль за функционированием скоринговых моделей.
	\item Защита данных и конфеденциальность (ст. 10) -- при эксплуатировании скоринговых моделей на основе ИИ, банк должен обеспечивать охрану и защиту персональных данных.
	\item Управление рисками систем ИИ (ст. 18) -- банки обязаны осуществлять непрерывный процесс выявления, анализа, оценки и минимизации рисков, связанных с использованием скоринговых моделей.
	\item Аудит систем ИИ (ст. 20) -- при проведении аудита скоринговых систем оценивается качество и правомерность использования данных для обучения моделей, а также проверку отсутствия в алгоритмах запрещенных функциональных возможностей, таких как дискриминация, манипуляция или использование уязвимостей заемщика (ст. 17).
	\item Библиотеки данных (ст. 27) -- банки обязаны использовать только законно полученные, качественные и актуальные данные для обучения скоринговых моделей с указанием сведений изготовителя по сбору данных (кредитная история, данные о доходах, платежах) в машиночитаемой форме, а также соблюдать заявленные цели использования данных.
\end{enumerate}

\section{Подходы к оценке и прогнозированию кредитного риска: классические модели и современные методы искусственного интеллекта}

Современные технологии ИИ находят применение в различных направлениях банковской деятельности. На основе анализа исследований были выделены 4 основных направления:
\begin{enumerate}[nosep]
	\item Чат-боты и робоконсультации. Предназначены для автоматизации консультаций и поддержки клиентов. Примером является Тинькофф-банк, который отмечает экономию 30-200 млн руб. в месяц за счет автоматизации, при этом более 40\% клиентов удовлетворяются ответом робота, а среднее время консультации сокращается примерно на 40 секунд\footnote{см. Там же. С. 15.}.
	\item Персонализация финансовых продуктов на основе анализа поведения клиентов. Рекомендательные модели Тинькофф-банка формируют персональные предложения, напоминая о покупках и подсказки по операциям в приложении\footnote{Бобкова Е.А. Применение искусственного интеллекта в банковской сфере: технологии и результаты // Выпускная квалификационная работа. Санкт-Петербургский государственный университет. 2025. С. 1--72.}.
	\item Антифрод-системы, предназначенные для выявления мошеннических трансакций в реальном времени. В Казахстане анти-фрод центр Национального банка за полгода зарегистрировал около 18 тыс. подзрительных трансакций и заблокировал порядка 1,5 млрд тг\footnote{\cite{NBRB} (дата обращения: 15.02.2026). Загл. с экрана.}, существенно снизив потенциальные потери клиентов.
	\item Кредитный скоринг -- ключевое направление для исследования. Американская компания Zest AI, разрабатывающая скоринговые модели на основе ИИ, показывает, что финансовые учреждения, внедрившие ее скоринговую систему, отмечают снижение рисков на 20\%\footnote{\cite{ZestAI} (дата обращения: 15.02.2026). Загл. с экрана.}, при сохранении объемов кредитования. И, наоборот, при сохранении прежнего уровня риска доля одобрений может быть увеличина на 15-20\%\footnote{см. Там же. С. 17.}.
\end{enumerate}

Качество скоринговых моделей напрямую зависит от набора используемых признаков. Традиционно при оценке кредитоспособности заемщиков используются такие классические признаки, как кредитнаяя история заемщика (наличие и частота просроченных платежей), уровень дохода, стаж занятости, показатель долговой нагрузки, возраст, пол и семейное положение. Данные характеристики являются основой классических скоринговых моделей, и формируются на основе предоставленной информации из Первого кредитного бюро. 

С развитием технологий машинного обучения и появлением новых источников данных стали использоваться альтернативные признаки, позволяющие дополнить традиционные методы оценки заемщиков. К таким признакам относятся данные об оплате коммунальных услуг и аренды жилья, данные о безналичных платежах и переводах, а также телеком-данные, подразумевающие использование мобильного интернета. Анализ таких данных позволяет более точно оценивать платежеспособность заемщиков, особенно в случаях, когда кредитная история ограничена или отсутствует. 

На основе анализа научной лиетратуры систематизированы ключевые различия между традиционными и современными подходами на основе ИИ к оценке вероятности дефолта. В таблице \ref{tab:compare_default_approaches} представлено сопоставление традиционного и современного подходов.

\begin{table}[H]
	\centering
	\caption{Сопоставление подходов к оценке вероятности дефолта}
	\label{tab:compare_default_approaches}
	\small
	\renewcommand{\arraystretch}{1.1}
	\begin{tabular}{|>{\RaggedRight\arraybackslash}p{0.30\textwidth}|
			>{\RaggedRight\arraybackslash}p{0.32\textwidth}|
			>{\RaggedRight\arraybackslash}p{0.32\textwidth}|}
		\hline
		\textbf{Критерий сравнения} & \textbf{Традиционный подход} & \textbf{Современный подход} \\ \hline
		
		\textbf{Тип моделей} 
		& Логистическая регрессия, деревья решений 
		& Случайный лес, градиентный бустинг, нейронные сети \\ \hline
		
		\textbf{Данные для оценки} 
		& Традиционные: кредитная история, доход, занятость.
		& Традиционные и альтернативные: аренда, коммунальные платежи, мобильные транзакции и т.д. \\ \hline
		
		\textbf{Охват клиентов} 
		& Клиенты с емкой кредитной историей. С тонкой или отсутствующей историей хуже оценивает 
		& Наиболее инклюзивны за счет альтернативных данных и моделей, которые могут выявлять сложные закономерности \\ \hline
		
		\textbf{Адаптивность} 
		& Статичны, хуже подстраиваются под изменения 
		& Динамичны, подстраиваются под обновление данных \\ \hline
		
		\textbf{Интерпретируемость} 
		& Проще объяснить решение 
		& Проблема «черного ящика», т.е. возникают сложности при интерпретации \\ \hline
		
		\textbf{Устойчивость при стресс-тестировании}
		& Наименее устойчивы при изменении экономической ситуации 
		& Наиболее устойчивы при изменении экономической ситуации \\ \hline
		
		\textbf{Ограничения} 
		& Недостаточно данных для оценки дефолта новых клиентов, слабо выявляют нелинейные связи 
		& Требования к высокой обеспеченности конфиденциальных данных, необходимость регулярного мониторинга моделей \\ \hline
		
	\end{tabular}
	
\end{table}

\begin{center}
	\begin{minipage}{\textwidth}
	\fontsize{9}{11}\selectfont
	\justifying
	Источник: составлено автором на основе работ Addy W.A., Ajayi-Nifise A.O., Bello B.G., Tula S.T., Odeyemi O., Falaiye T., Vakrani D.S., Padhye P.S., Gupta S.K., Roy J.K., Nerlekar V.S., Parashar N., Kothandapani H.P., Nuka
	T.F., Ogunola A.A.
\end{minipage}
\end{center}

Традиционные модели кредитного скоринга преимущественно используют ограниченный набор социальных и финансовых характеристик заемщиков, включая кредитную историю, уровень дохода и занятость. Основным преимуществом данного подхода является интерпретируемость результатов. 

Современные методы, основанные на алгоритмах МО и ГО, позволяют использовать больше признаков, включая альтернативные данные, такие как арендные, коммунальные платежи и мобильные транзакции. Использование таких признаков позволяет сделать шире охват клиентов и выявить более сложные зависимости в данных.

Тем не менее у современных моделей есть ряд ограничений. В частности, результаты таких моделей сложнее интерпретировать и они требуют регулярного мониторинга. Кроме того, их наиболее эффективно применять в работе с большими объемами данных, а также обеспечивать высокий уровень конфиденциальности информации.

Таким образом, современные модели ИИ позволяют расширять возможности оценки кредитного риска и позволяют повысить точность прогнозирования дефолта заемщиков. При этом эффективность скоринговых моделей напрямую зависит от качества исходных данных и их предварительной обработки.
%\begin{longtable}{|>{\RaggedRight\arraybackslash}p{0.23\textwidth}|
%		>{\RaggedRight\arraybackslash}p{0.25\textwidth}|
%		>{\RaggedRight\arraybackslash}p{0.27\textwidth}|
%		>{\RaggedRight\arraybackslash}p{0.18\textwidth}|}
%	\caption{Сопоставление подходов к оценке вероятности дефолта}
%	\label{tab:compare_default_approaches} \\
%	\hline
%	\textbf{Критерий сравнения} & \textbf{Традиционный подход} & \textbf{Современный подход} & \textbf{Источники} \\ \hline
%	\endfirsthead
%	
%	\hline
%	\textbf{Критерий сравнения} & \textbf{Традиционный подход} & \textbf{Современный подход} & \textbf{Источники} \\ \hline
%	\endhead
%	
%	\textbf{Тип моделей} &
%	Логистическая регрессия, деревья решений &
%	Случайный лес, градиентный бустинг, нейронные сети &
%	Addy W.A., Ajayi-Nifise A.O., Bello B.G., Tula S.T., Odeyemi O., Falaiye T., Vakrani D.S., Padhye P.S., Gupta S.K., Roy J.K., Nerlekar V.S., Parashar N., Даркенбаева Д.К., Зиятбековой Г.З., Алтыбай А., Мекебаева Н.О., Жамангарина Д.С. \\ \hline
%	
%	\textbf{Данные для оценки} &
%	Традиционные: кредитная история, доход, занятость и т.д. &
%	Традиционные и альтернативные: аренда, коммунальные платежи, мобильные транзакции и т.д. &
%	Addy W.A., Ajayi-Nifise A.O., Bello B.G., Tula S.T., Odeyemi O., Falaiye T., Kothandapani P., Hariharan H. (2022); Nuka T.F., Ogunola A.A. (2024). \\ \hline
%	
%	\textbf{Охват клиентов} &
%	Клиенты с емкой кредитной историей. С тонкой или отсутствующей историей хуже оценивает &
%	Наиболее инклюзивны за счет альтернативных данных и моделей, которые могут выявлять сложные закономерности &
%	Kothandapani P., Hariharan H. (2022); Nuka T.F., Ogunola A.A. (2024); Addy W.A., Ajayi-Nifise A.O., Bello B.G., Tula S.T., Odeyemi O., Falaiye T. (2022). \\ \hline
%	
%	\textbf{Адаптивность} &
%	Статичны, хуже подстраиваются под изменения &
%	Динамичны, подстраиваются под обновление данных &
%	Addy W.A., Ajayi-Nifise A.O., Bello B.G., Tula S.T., Odeyemi O., Falaiye T. (2022); Vakrani D.S., Padhye P.S., Gupta S.K., Roy J.K., Nerlekar V.S., Parashar N. (2026). \\ \hline
%	
%	\textbf{Интерпретируемость} &
%	Проще объяснить решение &
%	Проблема «черного ящика», т.е. возникают сложности при интерпретации &
%	Vakrani D.S., Padhye P.S., Gupta S.K., Roy J.K., Nerlekar V.S., Parashar N. (2026); Nuka T.F., Ogunola A.A. (2024). \\ \hline
%	
%	\textbf{Устойчивость при стресс-тестировании} &
%	Наименее устойчивы при изменении экономической ситуации &
%	Наиболее устойчивы при изменении экономической ситуации &
%	Vakrani D.S. и соавт. (2026). \\ \hline
%	
%	\textbf{Ограничения} &
%	Недостаточно данных для оценки дефолта новых клиентов, слабо выявляют нелинейные связи &
%	Требования к высокой обеспеченности конфиденциальных данных, необходимость регулярного мониторинга моделей &
%	Nuka T.F., Ogunola A.A. (2024); Addy W.A. и соавт. (2024). \\ \hline
%	
%\end{longtable}




